\section{Keyword Index}
\label{sec:keyword-index}

This section indexes the keywords accepted by the unified FEMaster
input language. The compatibility column distinguishes supported
Abaqus-compatible syntax from FEMaster extensions and legacy aliases;
it does not imply that every option offered by another solver is
implemented. Detailed keyword pages explain supported alternatives.

Command names are displayed in the conventional spaced form used in the manual, for example \cmd{END PART}, \cmd{SOLID SECTION}, and \cmd{STEADY STATE DYNAMICS}. \cmd{INERTIALOAD} is shown with the exact native spelling currently accepted by FEMaster.

\small
\begin{longtable}{@{}p{0.27\linewidth}p{0.19\linewidth}p{0.46\linewidth}@{}}
\caption{Input keywords and their canonical documentation locations.}\label{tab:keyword-index}\\
\toprule
\textbf{Keyword} & \textbf{Compatibility} & \textbf{Canonical topic} \\
\midrule
\endfirsthead
\toprule
\textbf{Keyword} & \textbf{Compatibility} & \textbf{Canonical topic} \\
\midrule
\endhead

\cmd{HEADING} & Abaqus-compatible & Input Structure \\
\cmd{MODEL} & FEMaster extension & Input Structure; optional native root marker \\
\cmd{PART} & Abaqus-compatible & Model Definition \\
\cmd{END PART} & Abaqus-compatible & Model Definition \\
\cmd{ASSEMBLY} & Abaqus-compatible & Model Definition \\
\cmd{END ASSEMBLY} & Abaqus-compatible & Model Definition \\
\cmd{INSTANCE} & Abaqus-compatible & Model Definition \\
\cmd{END INSTANCE} & Abaqus-compatible & Model Definition \\
\cmd{NODE} & Abaqus-compatible & Model Definition \\
\cmd{ELEMENT} & Compatible + extensions & Model Definition; \cmd{ELEMENT} formulations \\

\midrule
\cmd{NSET} & Abaqus-compatible & Model Definition \\
\cmd{ELSET} & Abaqus-compatible & Model Definition \\
\cmd{SURFACE} & Abaqus-compatible & Model Definition \\
\cmd{SFSET} & FEMaster extension & Model Definition \\
\cmd{PRETENSIONSECTION} & FEMaster extension & Bolt Pretension \\

\midrule
\cmd{MATERIAL} & Abaqus-compatible & Materials \\
\cmd{ELASTIC} & Abaqus-compatible & Materials \\
\cmd{PLASTIC} & Abaqus-compatible & Materials; isotropic J2 plasticity \\
\cmd{HYPERELASTIC} & Abaqus-compatible & Materials \\
\cmd{DENSITY} & Abaqus-compatible & Materials \\
\cmd{EXPANSION} & Abaqus-compatible & Materials \\
\cmd{THERMALEXPANSION} & Legacy alias & Materials; legacy alias of \cmd{EXPANSION} \\

\midrule
\cmd{PROFILE} & FEMaster extension & Sections and Element Properties \\
\cmd{SOLID SECTION} & Abaqus-compatible & Sections and Element Properties \\
\cmd{SHELL SECTION} & Abaqus-compatible & Sections and Element Properties \\
\cmd{BEAM SECTION} & FEMaster extension & Sections and Element Properties \\
\cmd{BEAM GENERAL SECTION} & Abaqus-compatible & Sections and Element Properties \\
\cmd{TRUSS SECTION} & FEMaster extension & Sections and Element Properties \\
\cmd{POINTMASS} & FEMaster extension & Sections and Element Properties; NSET-based concentrated properties \\
\cmd{MASS} & Abaqus-compatible & Sections and Element Properties; \val{MASS} point elements \\
\cmd{ROTARY INERTIA} & Abaqus-compatible & Sections and Element Properties; \val{ROTARYI} point elements \\
\cmd{SPRING} & Abaqus-compatible & Sections and Element Properties; linear \val{SPRING1} point elements \\

\midrule
\cmd{ORIENTATION} & Abaqus-compatible & Coordinate Systems and Orientations \\
\cmd{TRANSFORM} & Abaqus-compatible & Coordinate Systems and Orientations \\

\midrule
\cmd{FIELD} & FEMaster extension & Fields \\
\cmd{NORMAL} & FEMaster extension & Fields \\
\cmd{INITIALCONDITION} & FEMaster extension & Fields; initial nodal temperature and velocity \\
\cmd{INITIALCONDITIONS} & Alias & Fields; same syntax as \cmd{INITIALCONDITION} \\

\midrule
\cmd{AMPLITUDE} & Abaqus-compatible & Loads and Boundary Conditions \\
\cmd{SUPPORT} & FEMaster extension & Loads and Boundary Conditions \\
\cmd{BOUNDARY} & Abaqus-compatible & Loads and Boundary Conditions \\
\cmd{PRETENSION} & FEMaster extension & Bolt Pretension \\
\cmd{CLOAD} & Compatible + extensions & Loads and Boundary Conditions \\
\cmd{DLOAD} & Compatible + extensions & Loads and Boundary Conditions; Abaqus labels and FEMaster vector form \\
\cmd{PLOAD} & FEMaster extension & Loads and Boundary Conditions \\
\cmd{DSLOAD} & Abaqus-compatible & Loads and Boundary Conditions \\
\cmd{VLOAD} & FEMaster extension & Loads and Boundary Conditions \\
\cmd{INERTIALOAD} & FEMaster extension & Loads and Boundary Conditions \\
\cmd{TEMPERATURE} & FEMaster extension & Loads and Boundary Conditions \\
\cmd{HEATFLUX} & FEMaster extension & Loads and Boundary Conditions \\
\cmd{CONVECTION} & FEMaster extension & Loads and Boundary Conditions \\

\midrule
\cmd{COUPLING} & Compatible + extensions & Constraints and Contact \\
\cmd{KINEMATIC} & Abaqus-compatible & Constraints and Contact; child of scoped \cmd{COUPLING} \\
\cmd{DISTRIBUTING} & Abaqus-compatible & Constraints and Contact; child of scoped \cmd{COUPLING} \\
\cmd{EQUATION} & Abaqus-compatible & Constraints and Contact \\
\cmd{MPC} & FEMaster extension & Constraints and Contact; two-node kinematics \\
\cmd{TIE} & FEMaster extension & Constraints and Contact \\
\cmd{RIGID BODY} & CalculiX-compatible & Constraints and Contact \\
\cmd{CONNECTOR} & FEMaster extension & Constraints and Contact \\
\cmd{RBM} & FEMaster extension & Constraints and Contact \\
\cmd{CONTACT} & FEMaster extension & Constraints and Contact \\

\midrule
\cmd{LOADCASE} & FEMaster extension & Analysis Procedures \\
\cmd{SUPPORTS} & FEMaster extension & Analysis Procedures \\
\cmd{LOADS} & FEMaster extension & Analysis Procedures \\
\cmd{END} & FEMaster extension & Analysis Procedures \\
\cmd{STEP} & Abaqus-compatible & Analysis Procedures \\
\cmd{STATIC} & Abaqus-compatible & Analysis Procedures \\
\cmd{FREQUENCY} & Abaqus-compatible & Analysis Procedures \\
\cmd{BUCKLE} & Abaqus-compatible & Analysis Procedures \\
\cmd{DYNAMIC} & Abaqus-compatible & Analysis Procedures \\
\cmd{STEADY STATE DYNAMICS} & Abaqus-compatible & Analysis Procedures \\
\cmd{END STEP} & Abaqus-compatible & Analysis Procedures \\

\midrule
\cmd{SOLVER} & FEMaster extension & Solver and Numerical Controls; in Abaqus \cmd{STEP} after the procedure \\
\cmd{CONSTRAINTMETHOD} & FEMaster extension & Solver and Numerical Controls \\
\cmd{NONLINEAR} & FEMaster extension & Solver and Numerical Controls \\
\cmd{TIME} & FEMaster extension & Solver and Numerical Controls \\
\cmd{NEWMARK} & FEMaster extension & Solver and Numerical Controls \\
\cmd{DAMPING} & FEMaster extension & Solver and Numerical Controls \\
\cmd{FREQUENCIES} & FEMaster extension & Solver and Numerical Controls \\
\cmd{NUMEIGENVALUES} & FEMaster extension & Solver and Numerical Controls \\
\cmd{SIGMA} & FEMaster extension & Solver and Numerical Controls \\
\cmd{WRITEEVERY} & FEMaster extension & Solver and Numerical Controls \\
\cmd{INITIALVELOCITY} & FEMaster extension & Solver and Numerical Controls \\
\cmd{INERTIARELIEF} & FEMaster extension & Solver and Numerical Controls \\
\cmd{REBALANCELOADS} & FEMaster extension & Solver and Numerical Controls \\

\midrule
\cmd{OVERVIEW} & FEMaster extension & Output and Diagnostics \\
\cmd{REQUESTSTIFFNESS} & FEMaster extension & Output and Diagnostics \\
\cmd{REQUESTSTGEOM} & FEMaster extension & Output and Diagnostics \\
\cmd{CONSTRAINTSUMMARY} & FEMaster extension & Output and Diagnostics \\
\cmd{OUTPUT} & Abaqus-compatible & Output; \key{FIELD} container accepted, selection is made by child node/element output cards \\
\cmd{NODE OUTPUT} & Abaqus-compatible & Output; selects nodal result Fields for the active step \\
\cmd{ELEMENT OUTPUT} & Abaqus-compatible & Output; selects element-derived result Fields for the active step \\
\cmd{NODE FILE} & CalculiX-compatible & Output; CalculiX-style alias selecting nodal result Fields \\
\cmd{EL FILE} & CalculiX-compatible & Output; CalculiX-style alias selecting element-derived result Fields \\
\cmd{EL OUTPUT} & CalculiX-compatible & Output; element-derived frame Fields \\
\cmd{NODE PRINT} & Accepted, ignored & Output; accepted and ignored \\
\cmd{EL PRINT} & Accepted, ignored & Output; accepted and ignored \\

\midrule
\cmd{TOPODENSITY} & FEMaster extension & Topology-weighted Linear Statics \\
\cmd{TOPOORIENT} & FEMaster extension & Topology-weighted Linear Statics \\
\cmd{TOPOEXPONENT} & FEMaster extension & Topology-weighted Linear Statics \\

\bottomrule
\end{longtable}
\normalsize

\subsection{Supported element formulations}

The complete list of supported element names, Abaqus-compatible
aliases, node connectivity, numerical formulations, and limitations
is provided directly under \cmd{ELEMENT} in
Section~\ref{sec:model-definition}. The element-type table and
individual descriptions are maintained there as the canonical reference.

\subsection{Coverage notes}

Abaqus-compatible and CalculiX-compatible describe supported
input spellings and data subsets, not full behavioral equivalence to
those solvers. For example, \cmd{DLOAD} accepts body accelerations,
force-density components, element-face pressure labels, and a
FEMaster-specific vector surface-traction form. The detailed
keyword page specifies the valid data layouts.

\cmd{COUPLING}, \cmd{KINEMATIC}, and \cmd{DISTRIBUTING}
support Abaqus-compatible scoped coupling definitions with
inclusive degree-of-freedom ranges. An additional FEMaster
direct form uses \key{TYPE=KINEMATIC} or \key{TYPE=STRUCTURAL}
and six selector values. Other Abaqus-specific coupling options
are outside the supported subset.

\cmd{EQUATION} and concentrated-element properties have
Abaqus-compatible input forms. The supported point-element
subset consists of isotropic \val{MASS}, diagonal \val{ROTARYI},
and constant linear point-to-ground \val{SPRING1}. Off-diagonal
rotary inertia, oriented point properties, \val{SPRING2}/\val{SPRINGA},
nonlinear springs, and frequency- or field-dependent spring data
are outside the supported subset. The FEMaster extension
\cmd{POINTMASS} combines diagonal properties directly on an NSET.

The optional native \cmd{MODEL} marker is documented as input-file framing rather than as a separate level in the model-definition hierarchy.

The \cmd{PLASTIC} material card is available with Abaqus-compatible syntax, with the same tabulated yield-stress/equivalent-plastic-strain rows. \cmd{SOLVER} is a FEMaster extension accepted inside
\cmd{STEP} after the active procedure has been declared.
Other numerical-control keywords have the scopes and restrictions
documented on their individual reference pages.

Field-output cards \cmd{NODE OUTPUT}, \cmd{ELEMENT OUTPUT}, \cmd{NODE FILE}, \cmd{EL FILE}, and \cmd{EL OUTPUT} actively configure the per-step frame-result selection. The first explicit request replaces the step's frame-output defaults. Compact analysis metadata such as eigenvalues, frequencies, modal participation, and buckling factors is always written when the corresponding analysis provides it. \cmd{OUTPUT, FIELD} is accepted as a compatibility container only, while \cmd{NODE PRINT} and \cmd{EL PRINT} are accepted and ignored.

The index, detailed thematic documentation, and applicable complete examples together describe the supported input keywords and element formulations. This keeps the manual aligned with the syntax users can actually run.
